\clearpage
\begin{examplebox}
\textbf{\textcolor{CExample}{例1｜小学加法就能理解：两点为什么对称？}}
\begin{center}\begin{tikzpicture}[font=\small,>=stealth,line join=round]
\draw[gray!40,rounded corners=3pt] (0,0) rectangle (14.2,8.3);\node at (7.1,7.700000000000001) {先看一对点：中心就是连线的中点};
\begin{scope}[shift={(1.3,1.5)},x=1.4cm,y=1.1cm]\draw[->,gray] (0,0)--(4.7,0) node[right]{$x$};\draw[->,gray] (0,0)--(0,4.8) node[above]{$y$};\draw[blue,thick] (1,2)--(3,4);\fill[orange!85!black] (1,2) circle(2.4pt);\fill[orange!85!black] (3,4) circle(2.4pt);\fill[violet] (2,3) circle(2.5pt);\node[below right=7pt] at (1,2){$Q(1,2)$};\node[right=7pt] at (3,4){$P(3,4)$};\node[below right=8pt,violet] at (2,3){$C(2,3)$};\end{scope}\node[anchor=west,align=left] at (8.4,4.5){横坐标平均为2\\[14pt]纵坐标平均为3\\[14pt]$\overrightarrow{CP}=-\overrightarrow{CQ}$};
\node[text width=12.6cm,align=center] at (7.1,.65) {不指定整条函数曲线；一对点的对称，不能证明整张图像对称。};
\end{tikzpicture}\end{center}



\textbf{题目：}点 $P(3,4)$ 与 $Q(1,2)$ 关于哪个点对称？若它们是函数图像上的点，
这两个点对应的函数值有什么关系？

\textbf{解：}想象用细线连起 $P,Q$，找到线段的正中间：
\[
 C=\left(\frac{3+1}{2},\frac{4+2}{2}\right)=(2,3).
\]
因而两点关于 $C(2,3)$ 对称。如果 $f(3)=4, f(1)=2$，就有
$f(3)+f(1)=6=2\times3$。
\textbf{注意：}这只能说明\textbf{这一对点}对称，不能凭两点就认定\textbf{整张图}对称；
需要检查全部图像上的点。

\medskip

\end{examplebox}
\clearpage
\begin{examplebox}
\textbf{\textcolor{CExample}{例2｜基础判断：直接读出对称中心}}
\begin{center}\begin{tikzpicture}[font=\small,>=stealth,line join=round]
\draw[gray!40,rounded corners=3pt] (0,0) rectangle (14.2,8.3);\node at (7.1,7.700000000000001) {左右点不必同高，但连线的中点始终是中心};
\begin{scope}[shift={(1.491304347826087,2.4166666666666665)},x=1.3043478260869565cm,y=0.5111111111111111cm]\draw[->,gray] (-0.3,0)--(4.3,0) node[right]{$x$};\draw[->,gray] (0,-1.5)--(0,7.5) node[above]{$y$};\draw[blue,thick,solid,domain=0.4:3.6,samples=180] plot(\x,{(\x-2)^3+3});\fill[orange!85!black] (1,2) circle(2.2pt);\fill[orange!85!black] (3,4) circle(2.2pt);\draw[orange!85!black] (1,2)--(3,4);\draw[gray,dashed] (2,-1.2)--(2,7.2);\draw[gray,dashed] (-.2,3)--(4.2,3);\fill[violet] (2,3) circle(2.6pt);\end{scope}\node[anchor=west,align=left,text width=5.0cm] at (8.4,4.0) {$f=(x-2)^3+3$\\[9pt]$P(3,4),\ Q(1,2)$\\[9pt]$\frac{3+1}{2}=2,\quad\frac{4+2}{2}=3$\\[9pt]中点 $C(2,3)$\\[9pt]$f(2+t)+f(2-t)=6$};
\node[text width=12.6cm,align=center] at (7.1,.65) {这条具体曲线示范一般公式；严格判定仍要求所有点都有对应点。};
\end{tikzpicture}\end{center}



\textbf{题目：} $f$ 的定义域为 $\mathbb R$，任意实数 $t$ 均有
\[
 f(2+t)+f(2-t)=6.
\]
求图像的对称中心，并求 $f(2)$。

\textbf{解：}两个输入以 $2$ 为中间数，两个输出的和为 $6$，所以中间高度为 $6/2=3$。
定义域是全体实数，成对条件满足，因此\textbf{对称中心为 $(2,3)$}。
令 $t=0$：$f(2)+f(2)=6$，故\textbf{$f(2)=3$}。

\medskip

\end{examplebox}
\clearpage
\begin{examplebox}
\textbf{\textcolor{CExample}{例3｜看起来没有“$a\pm t$”，照样一秒识别}}
\begin{center}\begin{tikzpicture}[font=\small,>=stealth,line join=round]
\draw[gray!40,rounded corners=3pt] (0,0) rectangle (14.2,8.3);\node at (7.1,7.700000000000001) {例3：先找平均高度，再补另一侧的函数值};
\begin{scope}[shift={(1,1.8)},x=1cm,y=.65cm]\draw[->,gray] (0,0)--(7,0) node[right]{$x$};\draw[->,gray] (0,0)--(0,8);\draw[blue,thick] (1,7)--(5,3);\foreach \X/\Y in {1/7,5/3}{\fill[orange!85!black] (\X,\Y) circle(2.4pt);}\fill[violet] (3,5) circle(2.6pt);\node[above right=7pt] at (1,7){$(1,7)$};\node[right=7pt] at (5,3){$(5,3)$};\node[below left=7pt,violet] at (3,5){$(3,5)$};\end{scope}\node[anchor=west,align=left] at (9,4.5){两输入和为6\\[12pt]两输出和为10\\[12pt]平均得中心 $(3,5)$};
\node[text width=12.6cm,align=center] at (7.1,.65) {已知 $f(1)=7$，对应 $f(5)=10-7=3$；只画题目保证的一对点。};
\end{tikzpicture}\end{center}



\textbf{题目：}设 $f$ 的定义域为 $\mathbb R$，且
\[
 f(x)+f(6-x)=10.
\]
求对称中心；已知 $f(1)=7$，求 $f(5)$。

\textbf{解：} $x$ 与 $6-x$ 的平均数为 $3$，函数值的平均数为 $5$，
所以中心为\textbf{$(3,5)$}。取 $x=1$：
\[
 f(1)+f(5)=10\quad\Rightarrow\quad f(5)=10-7=\boxed{3}.
\]
\textbf{快捷记忆：}两个输入配对相加得 $6$；两个输出配对相加得 $10$。

\medskip

\end{examplebox}
\clearpage
\begin{examplebox}
\textbf{\textcolor{CExample}{例4｜不用解析式：一串函数值配对求和（高考常见思路）}}
\begin{center}\begin{tikzpicture}[font=\small,>=stealth,line join=round]
\draw[gray!40,rounded corners=3pt] (0,0) rectangle (14.2,8.3);\node at (7.1,7.700000000000001) {例4：每一对平均高度为5，中间那项也是5};
\begin{scope}[shift={(1.5285714285714287,1.8921052631578945)},x=0.8571428571428571cm,y=0.4842105263157894cm]\draw[->,gray] (-0.5,0)--(6.5,0) node[right]{$x$};\draw[->,gray] (0,-0.5)--(0,9) node[above]{$y$};\draw[blue,thick,solid,domain=0:6,samples=180] plot(\x,{\x+2});\fill[orange!85!black] (0,2) circle(2.2pt);\fill[orange!85!black] (1,3) circle(2.2pt);\fill[orange!85!black] (2,4) circle(2.2pt);\fill[orange!85!black] (3,5) circle(2.2pt);\fill[orange!85!black] (4,6) circle(2.2pt);\fill[orange!85!black] (5,7) circle(2.2pt);\fill[orange!85!black] (6,8) circle(2.2pt);\end{scope}\node[anchor=west,align=left,text width=5.0cm] at (8.4,4.0) {图取示例 $f=x+2$\\[9pt]$f(0)+f(6)=2+8=10$\\[9pt]$f(1)+f(5)=3+7=10$\\[9pt]$f(2)+f(4)=4+6=10$\\[9pt]$f(3)=5$，总和 $35$};
\node[text width=12.6cm,align=center] at (7.1,.65) {具体值用于看清配对；一般题不必知道各项值，只需每一对的和为10。};
\end{tikzpicture}\end{center}



\textbf{题目：} $f$ 定义在 $\mathbb R$ 上，满足 $f(x)+f(6-x)=10$。
求 $S=f(0)+f(1)+\cdots+f(6)$。

\textbf{解：}配对而不是逐个求值：
\[
\begin{array}{rcl}
 f(0)+f(6)&=&10,\\
 f(1)+f(5)&=&10,\\
 f(2)+f(4)&=&10.
\end{array}
\]
剩下中间那项 $f(3)$，令 $x=3$ 得 $2f(3)=10$，所以 $f(3)=5$。
于是 $S=3\times10+5=\boxed{35}$。

\textbf{举一反三：}若以 $a$ 为中间数，左右各取 $n$ 个对称位置，
且这些位置和 $a$ 都在定义域内，则 $2n+1$ 个函数值的和等于 $(2n+1)b$。
原因是每一对的平均值都是 $b$，中间那项也是 $b$。

\medskip

\end{examplebox}
\clearpage
\begin{examplebox}
\textbf{\textcolor{CExample}{例5｜只知道右边的解析式，怎样算左边？}}
\begin{center}\begin{tikzpicture}[font=\small,>=stealth,line join=round]
\draw[gray!40,rounded corners=3pt] (0,0) rectangle (14.2,8.3);\node at (7.1,7.700000000000001) {例5：点对称允许跳跃，中心单独取值};
\begin{scope}[shift={(3.242857142857143,2.092307692307692)},x=0.8571428571428571cm,y=0.8846153846153845cm]\draw[->,gray] (-2.5,0)--(4.5,0) node[right]{$x$};\draw[->,gray] (0,-0.5)--(0,4.7) node[above]{$y$};\draw[violet,thick,dashed,domain=-2:0.999,samples=180] plot(\x,{4-ln(2-\x)/ln(2)});\draw[blue,thick,solid,domain=1.001:4,samples=180] plot(\x,{ln(\x)/ln(2)});\fill[orange!85!black] (1,2) circle(2.2pt);\draw[blue,fill=white] (1,0) circle(2.8pt);\draw[blue,fill=white] (1,4) circle(2.8pt);\end{scope}\node[anchor=west,align=left,text width=5.0cm] at (8.4,4.0) {中心 $(1,2)$\\[9pt]右：$x>1,\ f=\log_2x$\\[9pt]左：$x<1$\\[4pt]$f=4-\log_2(2-x)$\\[9pt]中心点 $f(1)=2$\\[9pt]无需连续};
\node[text width=12.6cm,align=center] at (7.1,.65) {空心点 $(1,0),(1,4)$ 不取；中间实心点 $(1,2)$ 单独定义。};
\end{tikzpicture}\end{center}



\textbf{题目：} $f$ 的定义域为 $\mathbb R$，且
$f(1+t)+f(1-t)=4$；当 $x>1$ 时，$f(x)=\log_2 x$。
求 $f(-1)$，再写出 $x<1$ 时的 $f(x)$。

\textbf{解：} $-1$ 关于 $x=1$ 的对称位置是 $3$，因为 $-1+3=2$。
因此
\[
 f(-1)=4-f(3)=\boxed{4-\log_2 3}.
\]
对任意 $x<1$，它的对称位置 $2-x>1$，可以代入已知解析式：
\[
\boxed{\ f(x)=4-f(2-x)=4-\log_2(2-x),\quad x<1.\ }
\]
另外令 $t=0$ 得 $f(1)=2$。\textbf{提醒：}本题并没有说连续，
因此 $f(1)=2$ 与 $x>1$ 的表达式在 $x\to1^+$ 时趋近 $0$ 并不矛盾。

\medskip

\end{examplebox}
\clearpage
\begin{examplebox}
\textbf{\textcolor{CExample}{例6｜别把“右边递增”误判成“左边递减”}}
\begin{center}\begin{tikzpicture}[font=\small,>=stealth,line join=round]
\draw[gray!40,rounded corners=3pt] (0,0) rectangle (14.2,8.3);\node at (7.1,7.700000000000001) {横纵顺序同时反转，中心对称两侧同向单调};
\begin{scope}[shift={(1.2428571428571429,2.8)},x=1.4285714285714288cm,y=0.575cm]\draw[->,gray] (-0.1,0)--(4.1,0) node[right]{$x$};\draw[->,gray] (0,-2)--(0,6) node[above]{$y$};\draw[blue,thick,solid,domain=0.5:3.5,samples=180] plot(\x,{(\x-2)^3+2});\fill[orange!85!black] (0.5,-1.375) circle(2.2pt);\fill[orange!85!black] (1,1) circle(2.2pt);\fill[orange!85!black] (3,3) circle(2.2pt);\fill[orange!85!black] (3.5,5.375) circle(2.2pt);\end{scope}\node[anchor=west,align=left,text width=5.0cm] at (8.4,4.0) {示例 $f=(x-2)^3+2$\\[9pt]中心 $(2,2)$\\[9pt]左侧 $0.5<1$\\[9pt]右侧对应 $3.5>3$\\[9pt]$f(0.5)<f(1)$ 仍成立};
\node[text width=12.6cm,align=center] at (7.1,.65) {图取例6的一个具体函数；证明用比较式，适用于所有满足题设的函数。};
\end{tikzpicture}\end{center}



\textbf{题目：} $f$ 定义在 $\mathbb R$ 上，$f(2+t)+f(2-t)=4$，
并且在 $[2,+\infty)$ 上严格递增。证明它在 $(-\infty,2]$ 上也严格递增。

\textbf{解：}由题设，$f(x)=4-f(4-x)$。任取 $x_1<x_2\le2$，则
\[
4-x_1>4-x_2\ge2.
\]
已知右侧递增，所以 $f(4-x_1)>f(4-x_2)$。两边都用 $4$ 去减，
不等号转向：
\[
4-f(4-x_1)<4-f(4-x_2),
\]
也就是 $f(x_1)<f(x_2)$，证明完毕。

\textbf{直觉解释：}从右侧走向中心时函数值在下降；转半圈到左侧，
横坐标与纵坐标的前后顺序\textbf{一起反过来}，于是仍为递增。

\medskip

\end{examplebox}
\clearpage
\begin{examplebox}
\textbf{\textcolor{CExample}{例7｜对称中心不在图像上，算不算中心对称？}}
\begin{center}\begin{tikzpicture}[font=\small,>=stealth,line join=round]
\draw[gray!40,rounded corners=3pt] (0,0) rectangle (14.2,8.3);\node at (7.1,7.700000000000001) {中心可以不在图上：两支曲线仍绕它成对};
\begin{scope}[shift={(2.1,2.4166666666666665)},x=1.0cm,y=0.5111111111111111cm]\draw[->,gray] (-1,0)--(5,0) node[right]{$x$};\draw[->,gray] (0,-1.5)--(0,7.5) node[above]{$y$};\draw[blue,thick,solid,domain=-0.5:1.7,samples=180] plot(\x,{1/(\x-2)+3});\draw[blue,thick,solid,domain=2.3:4.5,samples=180] plot(\x,{1/(\x-2)+3});\fill[orange!85!black] (1,2) circle(2.2pt);\fill[orange!85!black] (3,4) circle(2.2pt);\draw[gray,dashed] (2,-1.4)--(2,7.4);\draw[gray,dashed] (-.9,3)--(4.9,3);\draw[violet,fill=white] (2,3) circle(3pt);\end{scope}\node[anchor=west,align=left,text width=5.0cm] at (8.4,4.0) {$f=\frac1{x-2}+3$\\[9pt]中心 $(2,3)$\\[9pt]$D=\mathbb R\setminus\{2\}$\\[9pt]$f(1)=2,\quad f(3)=4$\\[9pt]但 $f(2)$ 不存在};
\node[text width=12.6cm,align=center] at (7.1,.65) {紫色空心定位点不属于函数图像；两条虚线是渐近线，不是曲线。};
\end{tikzpicture}\end{center}



\textbf{题目：}求 $f(x)=\dfrac1{x-2}+3$ 的对称中心。

\textbf{解：}定义域 $D=\mathbb R\setminus\{2\}$，
若 $x\ne2$，则 $4-x\ne2$，所以定义域关于 $x=2$ 对称。
\[
\begin{aligned}
 f(x)+f(4-x)
 &=\left(\frac1{x-2}+3\right)
   +\left(\frac1{2-x}+3\right)\\
 &=6.
\end{aligned}
\]
故中心为\textbf{$(2,3)$}。不过 $f(2)$ 根本不存在，点 $(2,3)$ 并不在图像上。
\textbf{这就是检验“中心一定是曲线上的点”这一误解的反例。}

\medskip

\end{examplebox}
\clearpage
\begin{examplebox}
\textbf{\textcolor{CExample}{例8｜进阶联系：两个一样高的对称中心}}
\begin{center}\begin{tikzpicture}[font=\small,>=stealth,line join=round]
\draw[gray!40,rounded corners=3pt] (0,0) rectangle (14.2,8.3);\node at (7.1,7.700000000000001) {同一高度的两个中心，半转两次得到水平平移};
\begin{scope}[shift={(2.463636363636364,3.9499999999999997)},x=0.5454545454545454cm,y=1.5333333333333332cm]\draw[->,gray] (-2.5,0)--(8.5,0) node[right]{$x$};\draw[->,gray] (0,-1.5)--(0,1.5) node[above]{$y$};\draw[blue,thick,solid,domain=-2:8,samples=180] plot(\x,{sin(deg(pi*\x/3))});\fill[orange!85!black] (1,0.8660254037844386) circle(2.2pt);\fill[orange!85!black] (-1,-0.8660254037844386) circle(2.2pt);\fill[orange!85!black] (7,0.8660254037844386) circle(2.2pt);\end{scope}\node[anchor=west,align=left,text width=5.0cm] at (8.4,4.0) {图取 $f=\sin\frac{\pi x}{3}$\\[9pt]中心 $(0,0)$、$(3,0)$\\[9pt]$1\to-1\to7$\\[9pt]高度依次 $\frac{\sqrt3}2,-\frac{\sqrt3}2,\frac{\sqrt3}2$\\[9pt]相差 $6$，回到同高};
\node[text width=12.6cm,align=center] at (7.1,.65) {正文对任意函数的配对关系推出 $f(x+6)=f(x)$；6只保证是一个周期。};
\end{tikzpicture}\end{center}



\textbf{题目：}某函数定义域为 $\mathbb R$，图像分别关于 $(0,0)$、$(3,0)$ 中心对称。
证明 $6$ 是它的一个周期。

\textbf{解：}两次使用对称关系：
\[
 f(-x)=-f(x),\qquad f(6-x)=-f(x).
\]
将第二式中的 $x$ 换成 $-x$，有
$f(6+x)=-f(-x)=f(x)$。
因此\textbf{$f(x+6)=f(x)$，即 $6$ 是一个周期}。
注意不要求 $6$ 是\textbf{最小正周期}。
\end{examplebox}
